\section{Discussion}

SPINE now supports a residual-domain interpretation of the load-normalized spine-neck ratio. The ratio $\rhoload=\Rneck/\Rin$ encodes the classical local divider expectation $\Gamma_{\mathrm{div}}=1/(1+\rhoload)$; SMI remains only as legacy shorthand for this quantity in code-facing outputs. The scientific work begins after that analytic reference is stated: transient conductance-based simulations usually produced peak local transfer below the DC divider, and the residuals identify where a single scalar loses information.

This reframing is important because it separates a true but limited electrical expectation from broader prediction. The low-$\rhoload$ reference case was close to the divider limit, while intermediate and high reference cases were substantially below the divider in peak transfer. Across the full derived table, residuals were usually negative and sometimes large for the separate-peak metric. That pattern is not a failure to reproduce the divider; it is evidence that the divider is the low-frequency baseline and that peak synaptic responses live in a richer dynamic system.

The practical use pattern is therefore two-step. Use the load-normalized ratio to report first-pass local head-to-dendrite divider context. Then report residuals, component resistances, impedance, transfer gain, synaptic conductance scale, active state, morphology, and uncertainty whenever the target is peak transfer, head amplitude, somatic voltage, nonlinear response, class membership, or electrical equivalence. For experimental interpretation, that means reporting neck resistance together with dendritic/somatic load rather than geometry alone, and treating $\Gamma_{\mathrm{div}}$ as a transparent reference rather than as evidence that an observed transient response has been explained.

The ratio succeeds locally because it compares a series neck resistance with the load driven by that resistance. This is the classical cable-theoretic intuition for spine-neck isolation \citep{Rall1964,SegevRall1988,Yuste2013,Magee2000Review}. It also explains why the matched-neck control is meaningful: when $\Rneck$ is fixed, different dendritic loads produce different divider expectations.

The same logic explains why the ratio fails outside its target. Head amplitude depends on absolute synaptic conductance, driving-force reduction, local capacitance, and input impedance. Somatic transfer adds dendrite-to-soma filtering, a distinction related to location-dependent dendritic propagation and somatic EPSP normalization \citep{MageeCook2000}. Active state, NMDA voltage dependence, clustered input, and timing can change voltage trajectories without changing the morphology-derived ratio. Equal $\rhoload$ cannot imply electrical equivalence because many circuits share a ratio while differing in absolute RC constants, impedance phase, conductance drive, and morphology.

This does not make the ratio useless. It makes it conditional. A scalar descriptor can be valuable when its target is explicit and its residuals are measured. It becomes misleading only when the local divider coordinate is stretched into an amplitude, somatic, nonlinear, or categorical claim.

The radius audit highlights a measurement boundary. Because $\Rneck \propto 1/r_n^2$, modest radius error can produce large resistance and ratio error, and spine-neck geometry is itself plastic and difficult to measure precisely \citep{Tonnesen2014SpineNeckPlasticity,TonnesenNagerl2016}. The N=768 radius screen produced a 24.0\% design-fraction class flip, with all 11 intermediate-$\rhoload$ combinations crossing a class boundary. This supports continuous reporting with explicit uncertainty more strongly than thresholded low/intermediate/high labels.

The same N=768 ensemble contained no high-$\rhoload$ rows. This absence should be treated as a design limitation. The targeted high-$\rhoload$ extension added here is useful as a bounded stress test, but it is not a substitute for a full uncertainty ensemble spanning that regime. High-isolation behavior is therefore supported by reference cases, designed challenge rows, and targeted checks, whereas uncertainty stability is supported only for sampled low/intermediate combinations.

The from-scratch passive implementation is computationally credible within the scope tested. Internal unit, current-balance, limiting-case, timestep, impedance, and BE-versus-CN checks support numerical self-consistency. More importantly for external credibility, an independent direct-matrix implementation reproduced baseline passive traces and metrics to round-off, and DC analytic benchmarks checked the dendrite-soma load and divider limits.

Brian2 added a useful optional implementation-check anchor, but the scope remains narrow. Brian2 was installed in an isolated validation environment and the Brian2-hosted scoped checks reproduced representative passive baselines, one distributed-neck passive case, and one AMPA/NMDA synapse-only case within predefined tolerances. However, native Brian2 state updaters were not the final successful integration path for the stiff compact circuit, and restrained active channels were not reproduced in Brian2. The strongest supported statement is therefore narrower than global simulator validation: the passive reference equations are consistent with independent matrix/DC benchmarks and representative Brian2-hosted checks. Extending native cross-simulator checks to broader morphologies, active protocols, and uncertainty rows remains useful future work.

The results fit naturally with mixed experimental evidence about spine electrical isolation. Some optical and uncaging studies support substantial compartmentalization in particular contexts, whereas other work emphasizes modest isolation, indicator limits, and context-dependent interpretation \citep{Araya2006,Kwon2017,Cornejo2022,Popovic2015,Zecevic2023}. A manuscript that reports only a single ratio would flatten these differences. A manuscript that reports the ratio, divider expectation, residuals, and target-specific descriptors can connect morphology to physiology without overclaiming.

The summary-level external decision gate reinforces the reporting lesson without broadening the validation claim. The machine-recoverable Harnett and Popovic anchors suggest that some apparent high- versus low-isolation contrast can be reframed in load-normalized coordinates, but the evidence remains summary-level only: Harnett is a model/summary anchor rather than recovered paired measured per-spine data, Popovic measured and model-fit rows are kept separate, Kwon and Cornejo remain context/exclusion rows, and transfer proxies are not SPINE residuals. The overlay motivates reporting dendritic load together with neck resistance, but it is not a per-spine reanalysis, not biological validation, not a prevalence estimate, and not a controversy resolution. A true empirical test would require paired same-spine $\Rneck$, local dendritic load, and transfer definitions from the same preparation.

The current framework remains intentionally limited. The baseline soma is a lumped downstream load, passive morphologies are procedural rather than cell-specific reconstructions, active conductances are generic stress tests rather than calibrated densities, and uncertainty ensembles are deterministic sensitivity designs rather than posterior biological distributions. SPINE focuses on electrical behavior and does not model biochemical calcium reaction-diffusion, receptor trafficking, signaling cascades, or structural plasticity.

The most valuable next studies would combine voltage imaging, morphology measurements, and electrophysiology in preparations where neck-radius uncertainty can be propagated explicitly. Modeling extensions should prioritize calibrated cell-type active conductances, morphology-diverse posterior inference, analysis-defined high-$\rhoload$ uncertainty coverage, and fuller native cross-simulator checks. Accordingly, the strongest supported claim is that the divider is expected and the residuals are informative.

The load-normalized spine-neck ratio is therefore a compact coordinate for the classical local voltage divider, not a universal electrical phenotype. Its most defensible use is to state the expected low-frequency head-to-dendrite transfer and then measure the residual from that expectation. In SPINE, those residuals show when transient conductance dynamics, impedance, synaptic drive, active state, morphology, downstream filtering, and measurement uncertainty make the scalar descriptor incomplete.
